\scanpage{24}

qui contient $\widehat{\unsure{\Sigma|E}}$ comme sous-extension. On désigne par $\mathbb{\Gamma}$ le commutant de $\unsure{\mathfrak S_3|\mathcal S_3}$ dans $\widehat{\widehat\Gamma}$, \inserted{(qui contient $\dot\tau$)} et par $\mathbb E$ l'extension correspondante de $\mathbb{\Gamma}$ par $\widehat\pi$. Notons déjà que

\[
\left\{\begin{aligned}
\mathbb{\Gamma}\cap\Gamma\deleted{\unsure{\mathfrak S_3}}&=\{1,\tau\}\\
\mathbb{\Gamma}\cap\mathfrak S_3&=\{1\}
\end{aligned}\right.\qquad\deleted{\unsure{\mathrm{NB}}}
\]

% slanted
\margin{\uline{NB} \diagram{$\begin{array}{ccc}&\mathfrak S_3\times\widehat{\widehat\Gamma}^{!}&\\&\rotatebox{90}{$\subset$}&\\&\widehat{\widehat\Gamma}\supset\deleted{\unsure{\Gamma}}\mathfrak S_3\times\mathbb{\Gamma}&\end{array}$}
-- j'ignore si cette inclusion est une égalité, i.e. si $\mathbb{\Gamma}\subset\widehat{\widehat\Gamma}^{!}$ est une égalité}

On veut donner d'abord une description plus symétrique des éléments de $\widehat{\widehat\Sigma}^{!}$, sans faire jouer un rôle particulier à $L_0$ engendré par $l_0$ $(L_0=l_0^{\widehat{\mathbb Z}})$. Un $u\in\unsure{\widehat\Sigma|\widehat{\widehat\Sigma}}$ est \inserted{peut être} défini par un multiplicateur

\[
\unsure{p|\mathfrak p}=\chi(u)\in\widehat{\mathbb Z}^{*} \tag*{(19)}
\]

et des

\[
g_i\in\widehat{\unsure{\deleted{G}\pi|\pi}}\qquad i=0,1,\infty \tag*{(20)}
\]

\unsure{par|pour}

\[
u(l_i)=\mathrm{int}(g_i^{\unsure{\deleted{-1}|-1}})\,l_i^p \tag*{(21)}
\]

les données (19), (20) étant liées par la relation

\[
\bigl(\mathrm{int}(g_\infty^{\unsure{\deleted{-1}|-1}})\,l_\infty^p\bigr)
\bigl(\mathrm{int}\,g_1^{\deleted{-1}}\cdot l_1^p\bigr)
\bigl(\mathrm{int}(g_0^{\deleted{-1}})\,l_0^p\bigr)=1 \tag*{(22)}
\]

et par le fait que les $u(l_i)$ définis par (21) engendrent $\widehat\pi$ topologiquement. Pour $u$ fixé\unsure{,} $p$
% archive sheet number: 23
